\documentclass[10pt,a4paper]{amsart}
\usepackage[margin=19mm]{geometry}
\usepackage{amsmath,amssymb,mathtools}
\usepackage[scr=rsfs,cal=euler]{mathalfa}
\usepackage{xfrac}
\usepackage[unicode,hidelinks,bookmarks=false]{hyperref}
\newcommand{\ie}{i.e.\ }
\newcommand{\cf}{cf.\ }
\newcommand{\resp}{respectively\ }
\newcommand{\Subsection}{\subsection}
\providecommand{\nobreakdash}{\nobreak\hbox{-}}
\setlength{\emergencystretch}{2em}
\multlinegap0pt
\usepackage{bm}
\usepackage{bbm}
\usepackage{dsfont}
\usepackage{xspace}
\usepackage{cancel}
\usepackage{mathtools}
\usepackage{amsthm}
\newtheorem{theorem}{Theorem}
\newtheorem{lemma}[theorem]{Lemma}
\title[The J-equation at the birational minimal slope]{The J-equation at the birational minimal slope}
\author{Junbang Liu}
\date{Source: arXiv:2609.28411v1 (CC BY 4.0)}
\begin{document}
\maketitle

\noindent\emph{Source and licence.} The J-equation at the birational minimal slope. Version arXiv:2609.28411v1 (CC BY 4.0), distributed under the Creative Commons Attribution 4.0 International licence, as recorded in the arXiv licence metadata for this exact version and shown on the abstract page https://arxiv.org/abs/2609.28411v1. Full rights notice: licenses/arxiv-2609.28411-CC-BY-4.0.txt.\par\smallskip

\noindent\textit{Editorial scope.} Counted unit: the Lemma whose source label is \texttt{lem:bergman-trace} in the article (source0.tex lines 920--931, source label \texttt{lem:bergman-trace}), delivered together with the article's own complete original proof of it (source0.tex lines 932--1014, one \texttt{proof} environment of 83 lines). No mathematics is written, repaired or re-derived: statement and proof are the article's own text line for line. Required context retained uncounted: none. Every definition, display and label of those lines is retained unchanged. Omitted from this package and not redistributed: 1--919, 1015--4035. Nothing inside a retained excerpt is omitted or altered: every retained body is delivered exactly as the article's lines are written, so no omission receipt of any kind accompanies this package. Citations: the retained excerpts carry no citation command at all, so no bibliography entry of the article is transcribed and no reference is redistributed. Reference adaptation: none was needed and none was made; every reference inside the delivered unit resolves inside it, so no unresolved reference or citation is produced. Printed slips kept literal: no printed word, symbol, hypothesis or number of the counted statement or of its proof was altered. Layout and class: the article's own class is \texttt{amsart} with packages including fontenc, inputenc, lmodern, amsmath, amssymb, mathtools, mathrsfs, microtype, enumitem, needspace, tikz, hyperref, biblatex; the wrapper is the lane's pinned \texttt{amsart} editorial wrapper at 10pt on A4, which restates the theorem-like environments the retained text needs and adds the article's own macro definitions that the retained text needs (none), with extra wrapper packages (bm, bbm, dsfont, xspace, cancel, mathtools). The delivered numbering is the unit's own. Assets: no retained span contains a figure environment, a graphics-inclusion command, a PDF-page inclusion, a table or a TikZ picture; no image file is copied into this package, the package declares no asset dependency and no image file is redistributed with it. Licence: version arXiv:2609.28411v1 of this article is distributed under the Creative Commons Attribution 4.0 International licence, as recorded in the arXiv licence metadata for this exact version and shown on the abstract page https://arxiv.org/abs/2609.28411v1. A full rights notice is delivered at licenses/arxiv-2609.28411-CC-BY-4.0.txt. Characters: every retained excerpt is pure ASCII, so the delivered engine, XeLaTeX with its default Latin Modern text font, reports no missing character.
\begin{lemma}\label{lem:bergman-trace}
Let $\Omega\Subset\mathbb C^n$ be bounded pseudoconvex, and let $u$ be plurisubharmonic near $\overline\Omega$. Fix a constant positive form $B$, $a>0$, and $1\le q\le n$. Suppose
\begin{equation}\label{eq:bergman-full}
 P_{q,B}(H_{u*\varrho_r})\le a
 \quad\text{on }\Omega\text{ for all sufficiently small }r>0.
\end{equation}
For the kernel of $e^{-mu}d\lambda$, the potential $b_m=m^{-1}\log K_m(z,z)$ is smooth on $\{K_m(z,z)>0\}$, and
\begin{equation}\label{eq:bergman-partial}
 H_{b_m}>0,\qquad P_{q,B}(H_{b_m})\le a.
\end{equation}
For a smooth positive density $\rho\,d\lambda$ defined near $\overline\Omega$, the bound is $a/(1-C/m)$ for all sufficiently large $m$.
\end{lemma}

\begin{proof}
A constant linear change of coordinates reduces $B$ to $I$. Its constant Jacobian only multiplies the kernel by a constant. First assume $u$ is smooth; letting $\delta\downarrow0$ in the
hypothesis gives $P_{q,I}(H_u)\leq a$.  For $K_m(w,w)>0$, set
\[
 s_i(z)=(\bar z_i-\bar w_i)k_w(z).
\]
Reproduction gives $\langle s_i,f\rangle_m
=\langle k_w,(z_i-w_i)f\rangle_m=0, \forall f\in\mathscr H_m$, i.e. $s_i\perp\mathscr H_m$, so $s_i$ is the minimal-norm solution of $\bar\partial s_i=k_w\,d\bar z_i$. H\"ormander's estimate gives
\begin{equation}
 \|s_i\|_m^2\le\frac1m\int_\Omega(H_u^{-1})_{i\bar i}|k_w|^2e^{-mu}d\lambda.
\end{equation}
To relate this estimate to $b_m$, choose an orthonormal basis
with $\sigma_1=k_w$. The reproducing property implies
$\sigma_\ell(w)=0$ for every $\ell\ge2$.
Differentiating $b_m=m^{-1}\log\sum_\ell|\sigma_\ell|^2$ at
$w$, and direct computation gives 
\begin{equation}
 m(b_m)_{i\bar j}(w)=
 \frac{\sum_{\ell\ge2}\partial_i\sigma_\ell(w)\overline{\partial_j\sigma_\ell(w)}}{K_m(w,w)}.
\end{equation}
Every $f\in\mathscr H_m$ with $f(w)=0$ expands in the
remaining basis functions $\sigma_\ell$ with $\ell\ge 2$. Cauchy--Schwarz inequality therefore gives
\[
|\partial_i f(w)|^2
\le m(b_m)_{i\bar i}(w)K_m(w,w)\,\|f\|_m^2.
\]
In unitary coordinates diagonalizing $H_{b_m}(w)$, apply this to $f=(z_i-w_i)k_w$. Since $\Omega$ is bounded, $f\in\mathscr H_m$, and
\[
 1\le m(b_m)_{i\bar i}(w)\|s_i\|_m^2.
\]
Thus $H_{b_m}(w)>0$. Let $J$ index its $q$ largest reciprocal eigenvalues. Then
\[
 P_{q,I}(H_{b_m}(w))
 \le m\sum_{i\in J}\|s_i\|_m^2
 \le\int_\Omega\sum_{i\in J}(H_u^{-1})_{i\bar i}|k_w|^2e^{-mu}d\lambda\le a.
\]
For general $u$, take $u_\delta=u*\varrho_\delta\downarrow u$
on the fixed domain $\Omega$, and denote the corresponding
kernel by $K_{m,\delta}$.
Since $e^{-mu_\delta}\uparrow e^{-mu}$, the extremal formula
shows that $K_{m,\delta}(w,w)$ decreases as
$\delta\downarrow0$, and
\[
K_{m,\delta}(w,w)\ge K_m(w,w).
\]
For fixed $w$, the normalized kernels
\[
\frac{K_{m,\delta}(z,w)}{\sqrt{K_{m,\delta}(w,w)}}
\]
form a normal family, by local upper bounds for $u_\delta$
and the submean inequality. Any subsequential limit $f$
satisfies, by Fatou's lemma,
\[
\|f\|_m\le1,\qquad
|f(w)|^2=\lim_{\delta\downarrow0}K_{m,\delta}(w,w).
\]
The extremal formula for $K_m$ gives the opposite inequality
for this limit. Hence $K_{m,\delta}(w,w)\downarrow K_m(w,w)$.
The bounds
\[
|K_{m,\delta}(z,w)|^2
\le K_{m,\delta}(z,z)K_{m,\delta}(w,w)
\]
give local normality of the two-variable kernels.
Their limits are determined by the diagonal, so polarization and
Cauchy estimates yield $K_{m,\delta}\to K_m$ in
$C^\infty_{\rm loc}(\Omega\times\Omega)$.
On $\{K_m(z,z)>0\}$, the logarithmic Hessians converge.
Closedness of $\{H>0:P_{q,I}(H)\le a\}$ proves the desired inequality.

Finally, the density is absorbed into the weight:
$e^{-mu_\delta}\rho=e^{-m(u_\delta-m^{-1}\log\rho)}$.
Since $H_{u_\delta}\ge I/a$ and $H_{\log\rho}$ is bounded on
$\overline\Omega$, there is $C$, independent of $m,\delta$, with
\[
H_{u_\delta}-\frac1mH_{\log\rho}
\ge(1-C/m)H_{u_\delta}>0
\]
for $m>C$.
Inversion bounds its partial inverse trace by $a/(1-C/m)$.
Apply the smooth result to $u_\delta-m^{-1}\log\rho$, then let
$\delta\downarrow0$ by the preceding kernel-convergence argument.
\end{proof}

\end{document}
